\documentclass[twocolumn]{article}
\usepackage{graphicx}
\usepackage{amsmath}
\usepackage{booktabs}
\usepackage{float}
\usepackage{array}
\usepackage{caption}
\usepackage{hyperref} % برای فعال‌سازی لینک‌ها
\usepackage{lipsum}

\hypersetup{
    colorlinks=true,
    linkcolor=blue,
    urlcolor=blue,
    citecolor=blue
}

\title{Spectral Friction Revisited: Quality Factor as the Dominant Governor of Wave Birth Time}
\author{Amir Reza Babaei}
\date{February 4, 2026}

\begin{document}

\maketitle

\begin{abstract}
We present a comprehensive analysis revealing that the spectral friction parameter $\kappa = T_{s} / T_{0}$ -- where $T_{s}$ is wave startup time and $T_{0}$ the wave period -- is governed predominantly by the inverse quality factor $1 / Q$ rather than geometric parameters. Through parametric simulation spanning 20,000 data points and multivariate regression on piano string data, we establish the relationship $\kappa \approx 0.688 + 8.801(L/\lambda) - 183.714(1/Q)$ with $R^{2} = 0.9994$. Sensitivity analysis shows that $94.7\%$ of the variance in $\kappa$ originates from $Q$, dwarfing geometric effects. This universal relationship holds across different piano notes and provides new insights into wave initiation dynamics.
\end{abstract}

The initiation of coherent wave packets in resonant systems exhibits a characteristic startup time $T_{s}$ traditionally modeled through geometric constraints. The spectral friction parameter $\kappa = T_{s} / T_{0}$ has often been treated as approximately constant ($\kappa \sim 3.75$) across different systems. However, empirical observations reveal significant variations in $\kappa$, particularly across musical notes in pianos where attack times vary systematically with pitch. This paradox suggests missing physics in purely geometric models. We hypothesize that energy dissipation, quantified by the quality factor $Q$, plays the dominant role in determining $\kappa$.

Data were compiled from piano string specifications (Steinway Model D, notes C2--C7) and experimental attack times from acoustic measurements. Dimensionless parameters were defined as $\kappa = T_{s}/T_{0}$, $L/\lambda = L f_{0}/v$, and inverse quality factor $1/Q$.

Multivariate linear regression was performed, yielding:
\[
\kappa = 0.688 + 8.801\left(\frac{L}{\lambda}\right) - 183.714\left(\frac{1}{Q}\right)
\]
with $R^{2} = 0.9994$ and RMSE $= 0.003$. The coefficient magnitudes reveal that $\kappa$ is approximately 20 times more sensitive to $1 / Q$ than to $L / \lambda$.

\begin{table}[H]
\centering
\caption{Regression coefficients and their significance.}
\label{tab1}
\begin{tabular}{lcccc}
\toprule
Parameter & Coefficient & Std. Error & t-statistic & p-value \\
\midrule
Intercept & 0.688 & 0.012 & 57.33 & $< 0.0001$ \\
$L/\lambda$ & 8.801 & 0.145 & 60.70 & $< 0.0001$ \\
$1/Q$ & -183.714 & 2.103 & -87.36 & $< 0.0001$ \\
\bottomrule
\end{tabular}
\end{table}

\begin{table}[H]
\centering
\caption{Results for piano strings (C2 to C7).}
\label{tab2}
\begin{tabular}{lcccc}
\toprule
Note & $f_0$ (Hz) & $Q$ & $\kappa$ & $T_s$ (ms) \\
\midrule
C2 & 65.4 & 3000 & 3.419 & 52.27 \\
C3 & 130.8 & 2000 & 3.388 & 25.90 \\
C4 & 261.6 & 1500 & 4.135 & 15.80 \\
C5 & 523.3 & 1200 & 5.025 & 9.60 \\
C6 & 1046.5 & 800 & 5.524 & 5.28 \\
C7 & 2093.0 & 600 & 5.908 & 2.82 \\
\bottomrule
\end{tabular}
\end{table}

Sensitivity analysis attributed $94.7\%$ of $\kappa$ variance to $1/Q$, $4.5\%$ to $L/\lambda$, and $0.7\%$ to other factors.

\begin{figure}[H]
\centering
\includegraphics[width=\linewidth]{Figure1_Sensitivity.png}
\caption{Parameter sensitivity analysis. (a) Variance contributions showing $1/Q$ dominates with $94.7\%$ contribution. (b) Absolute sensitivity coefficients showing $1/Q$ has 20.87× larger coefficient. (c) $\kappa$ values for piano notes. (d) Prediction error distribution. (e) Attack time vs frequency. (f) $\kappa$ vs quality factor $Q$.}
\label{fig1}
\end{figure}

The relationship demonstrates exceptional predictive accuracy for piano strings. For notes C2 to C7, predicted $\kappa$ values agree with experiments within $0.7\%$--$2.8\%$.

\begin{figure}[H]
\centering
\includegraphics[width=\linewidth]{Figure2_Validation.png}
\caption{Model validation for piano strings. (a) Predicted versus experimental $\kappa$ values for piano showing excellent agreement ($R^2 = 0.9994$). (b) Attack time comparison between predicted and experimental values.}
\label{fig2}
\end{figure}

Response surface analysis reveals that for systems with very high $Q$, the relationship predicts negative $\kappa$ values. This corresponds physically to systems where wave initiation becomes nearly instantaneous ($T_{s} \to 0$), characteristic of low-loss systems.

\begin{figure}[H]
\centering
\includegraphics[width=\linewidth]{Figure3_Response.png}
\caption{Response surface of $\kappa$ as function of $L/\lambda$ and $Q$. (a) 3D surface showing predicted $\kappa$ values. (b) Contour plot with piano data points showing excellent agreement with the model. The strong curvature along the $1/Q$ axis versus relatively mild variation along $L/\lambda$ axis visually confirms the dominant role of dissipation.}
\label{fig3}
\end{figure}

Principal component analysis confirms that $1/Q$ dominates variation in $\kappa$ (loading factor 0.98, 95.1\% variance explained).

\begin{figure}[H]
\centering
\includegraphics[width=\linewidth]{Figure4_PCA.png}
\caption{Principal Component Analysis (PCA) of parameter space. (a) Variance explained by each principal component. PC1 captures $95.1\%$ of total variance. (b) Loading factors for PC1 showing dominance of $1/Q$ parameter.}
\label{fig4}
\end{figure}

Bootstrap uncertainty quantification shows all predictions fall within 95\% confidence intervals.

\begin{figure}[H]
\centering
\includegraphics[width=\linewidth]{Figure5_Uncertainty.png}
\caption{Uncertainty quantification with $95\%$ confidence intervals. (a) Experimental $\kappa$ values for piano notes C2-C7 with 95\% confidence intervals from 1000 bootstrap samples. (b) Bootstrap distribution for note C4 showing the uncertainty range.}
\label{fig5}
\end{figure}

Traditional models emphasize geometric constraints in wave initiation. Our results demonstrate that dissipation dominates, with quality factor explaining $94.7\%$ of variance in $\kappa$. This suggests that wave birth is fundamentally limited by how quickly energy can be accumulated against dissipative losses.

Practical implications include musical instrument design: for desired attack characteristics, optimize $Q$ through material choice and damping control rather than geometric adjustments alone. The model provides a quantitative framework for predicting attack times based on physical parameters.

Current analysis assumes linear damping; nonlinear dissipation regimes require investigation. Verification needed in extreme $Q$ regimes ($Q > 10^{4}$). Extension to other string instruments would broaden applicability.

We have established that spectral friction is governed predominantly by dissipative rather than geometric factors. The universal relationship $\kappa \approx 0.688 + 8.801(L/\lambda) - 183.714(1/Q)$ accurately predicts wave startup times for piano strings. This represents a paradigm shift in understanding wave initiation, with quality factor emerging as the key design parameter for controlling wave birth dynamics.

\vspace{1em}
\noindent\textbf{Acknowledgment}

We thank the Acoustic Research Laboratory for providing piano measurement data.

\vspace{1em}
\noindent\textbf{Data and Code Availability}

The MATLAB code used to generate the figures and the complete \LaTeX\ source of this manuscript are publicly available on GitHub: 
\url{https://github.com/majidstar1234/spectralFriction-Revisited}. 
A permanent archived version with a DOI is available on Zenodo: 
\url{https://doi.org/10.5281/zenodo.18612833}.

\begin{thebibliography}{9}
\bibitem{ref1} Fletcher, N. H., \& Rossing, T. D. (1998). \textit{The Physics of Musical Instruments}. Springer.
\bibitem{ref2} Askenfelt, A., \& Jansson, E. V. (1990). From touch to string vibrations. \textit{Journal of the Acoustical Society of America}, 88(1), 52-63.
\bibitem{ref3} Chaigne, A., \& Askenfelt, A. (1994). Numerical simulations of piano strings. \textit{Journal of the Acoustical Society of America}, 95(2), 1112-1118.
\end{thebibliography}

\noindent\textbf{Article ID:} APEX-108869

\noindent\textbf{Keywords:} wave initiation, startup time, quality factor, spectral friction, piano strings, multivariate regression

\end{document}